Let $\mathcal R_\alpha(q)$ be the set of $k$ valid, unique database IDs
returned at budget $\alpha$ on a fixed index. Define
$\mathcal R_{\rm ref}(q)=\mathcal R_{\alpha_{\rm ref}}(q)$. The descending grid is
\begin{equation}
 A=(200,100,64,32,16,8,4,2),\qquad \alpha_{\rm ref}=A[0]=200.
 \label{eq:6-grid}
\end{equation}
For $S$ calibration queries $\mathcal Q_S$, define
\begin{equation}
 D_S(\alpha)=\sum_{q\in\mathcal Q_S}
 \left(k-|\mathcal R_\alpha(q)\cap\mathcal R_{\rm ref}(q)|\right),
 \qquad D_S(\alpha)\leq\epsilon.
 \label{eq:6-agreement}
\end{equation}
The evaluated $\epsilon=1$ permits one missing reference ID across the
whole sample. It ignores output order and distance differences; replacing
one ID counts once, although the symmetric difference contains two IDs.
The finite-grid objective is
\begin{equation}
 \widehat\alpha=\min\{\alpha\in A:D_S(\alpha)\leq\epsilon\}.
 \label{eq:6-selection-objective}
\end{equation}
The reference supplies a nonempty fallback. The objective minimises an
accepted multiplier; \cref{sec:6-3} states when bisection solves it exactly.

At $k=10$, the normalised allowance $\epsilon/(Sk)$ is 0.003125,
0.0015625 and 0.00078125 for $S=32,64,128$, and 0.0002 at $S=500$.
Increasing $S$ therefore changes both coverage and effective strictness.
Calibration and held-out evaluation use disjoint queries; subsamples
reuse a finite calibration pool (\cref{sec:7-5}).

\begin{figure}[htbp]
\centering
% Native schematic of the implemented decision and its offline evaluation.
\begin{tikzpicture}[
 font=\footnotesize\linespread{1}\selectfont,
 box/.style={draw=ThesisNavy!65,rounded corners=2pt,align=center,
 text width=2.8cm,minimum height=1.7cm,inner sep=4pt},
 arr/.style={-{Latex[length=2mm]},draw=ThesisNavy!80,thick}]
\node[box] (sample) at (0,0) {Calibration pool\\Even indices\\Sample $S$ queries};
\node[box,fill=ThesisNavy!7] (ref) at (3.6,0) {Reference search\\$\alpha_{\rm ref}=200$\\Top-$k$ IDs};
\node[box,fill=ThesisNavy!7] (test) at (7.2,0) {Bisection\\$D_S(\alpha)\leq1$\\Shared sample};
\node[box] (choice) at (10.8,0) {Chosen $\alpha$\\Caller applies it\\For batch reuse};
\node[box] (held) at (0,-3) {Held-out pool\\Odd query indices\\GT for evaluation};
\node[box,text width=6.4cm] (eval) at (7.2,-3)
 {Evaluate the selected arm\\Look up held-out recall and QPS\\from previously timed grid arms};
\draw[arr] (sample) -- (ref);
\draw[arr] (ref) -- (test);
\draw[arr] (test) -- (choice);
\draw[arr] (held) -- (eval);
\draw[arr,dashed] (choice.south) -- (10.8,-3) -- (eval.east);
\end{tikzpicture}

\caption{Workload-level calibration and its evaluation. Budget decisions
use reference IDs; GT assesses held-out quality. Selected-arm results
are looked up from timed fixed arms. Implementation boundaries are in
\cref{sec:a-budget-contract}.}
\label{fig:6-calibration}
\end{figure}

\subsection{What agreement controls}

Under \Cref{prop:6-monotonicity}, disagreement counts reference answers
excluded by primary screening:
\begin{equation}
 D_S(\alpha)=\sum_{q\in\mathcal Q_S}
       |\mathcal R_{\rm ref}(q)\setminus C_\alpha(q)|.
 \label{eq:6-reference-coverage}
\end{equation}
For valid fixed-size output sets, let $r_{\mathcal Q_S}$ denote GT recall
on these same queries. Each missing reference ID removes at most one GT hit, so
\begin{equation}
 r_{\mathcal Q_S}(\alpha)\geq r_{\mathcal Q_S}({\rm ref})
                         -\frac{D_S(\alpha)}{Sk}.
 \label{eq:6-relative-recall}
\end{equation}
This bound is sample-local and reference-relative. Agreement controls
additional loss relative to the reference; the reference's own quality
requires GT evaluation. \Cref{sec:7-recall-attribution} quantifies that
distinction with the reconstruction pilot.

\subsection{Sampling resolution and ordered testing}

Let $H_S$ count queries with any missing reference ID. Since $H_S\leq D_S$,
acceptance implies $H_S\leq1$. For a prespecified budget and i.i.d.
queries, the one-sided 95\% exact binomial upper limit after one event solves
\cite{nist_binomial}
\begin{equation}
 (1-u)^S+Su(1-u)^{S-1}=0.05.
 \label{eq:6-binomial-limit}
\end{equation}
The query-event limits are 0.03652 at $S=128$ and 0.009452 at 500.
Treating all $Sk$ result indicators as homogeneous independent Bernoulli
trials is stronger: the corresponding limits are 0.003701 and 0.0009484,
approximately $4.74/(Sk)$. Query i.i.d. alone does not justify dividing
the query-event limit by $k$; neither model certifies repeated subsamples
of this fixed pool.

For monotone disagreement and a fixed grid order, bisection selects the
accepted prefix. If its per-arm decisions are valid level tests, the
first-unsafe-arm argument gives fixed-sequence error control without a
Bonferroni factor \cite{learn_then_test}. The implemented one-slot rule
is a stability threshold; a GT-risk guarantee would require a valid test
for that target.

\paragraph{A finite-pool extension.}
Suppose every query has a fixed disagreement count $X_q(\alpha)$ under
the intended sample-size execution regime, unchanged by subset membership
or position. Let $N_{\rm cal}$ be the pool size and $n_j$ count queries with $X_q=j$.
Under uniform sampling without replacement, $D_S\leq1$ requires either
only zero-count queries or one single-count query \cite{nist_hypergeometric}:
\begin{equation}
 \Pr\{D_S(\alpha)\leq1\}
 =\frac{\binom{n_0}{S}+n_1\binom{n_0}{S-1}}{\binom{N_{\rm cal}}{S}}.
 \label{eq:6-finite-pool}
\end{equation}
Out-of-range binomial coefficients are zero. This is a conditional
acceptance model. Predicting selection risk additionally needs the
grid's acceptance structure and selected-arm quality map. The supplementary full-pool records reproduce every sampled curve by
summing the recorded $X_q$ over its actual sample IDs: 18,432 grid counts
match, including changes of execution path (\cref{sec:7-budget-trace}).
This validates the aggregate counts on those draws. Prediction for arbitrary
subsets still assumes stable per-query counts. Acceptance needs no labels;
predicting GT-loss crossings additionally uses the fixed-arm held-out
quality map.
